\section{Concluding remarks}\label{sec:conclusion}

Theorem~\ref{thm:main} shows that Foundation is unnecessary for the
equivalence between the existence of vector space bases and $\AC$,
settling the questions that appear in
\cites[p.~33]{Blass}{Karagila2017MO}[p.~423]{Morillon}[p.~121]{Philip2025}.

Blass~\cite[p.~33]{Blass} also asked whether, without Foundation, $\AC$
follows from the assertion that every linearly independent set in a
vector space can be extended to a basis. The following corollary
answers this affirmatively and also includes the assertion that every
spanning subset of a vector space contains a basis of that space,
studied by Halpern~\cite{Halpern1966}.

\begin{corollary}\label{cor:basis-principles}
In $\ZFm$, the following statements are equivalent:
\begin{enumerate}
\item $\AC$;
\item every vector space has a basis;
\item every vector space over a field of characteristic zero has a basis;
\item every linearly independent subset of a vector space is contained
in a basis of that space;
\item every spanning subset of a vector space contains a basis of
that space.
\end{enumerate}
\end{corollary}

\begin{proof}
Assume $(1)$. Given a linearly independent subset $I$ of a vector
space $V$, Zorn's lemma yields a maximal linearly independent subset
of $V$ containing $I$. Such a subset spans $V$ and is therefore a
basis, proving $(4)$. Similarly, given a spanning subset $S$ of $V$,
Zorn's lemma yields a maximal linearly independent subset $B$ of $S$.
Maximality implies that $S$ is contained in the span of $B$, so $B$
is a basis of $V$, proving $(5)$. These applications of Zorn's lemma
do not require Foundation.

Taking $I=\varnothing$ in $(4)$ gives $(2)$, as does taking $S=V$
in $(5)$. Clearly, $(2)$ implies $(3)$. Finally, $(3)$ implies that
every vector space over a field of characteristic zero
has a basis, so Theorem~\ref{thm:main} gives $(1)$.
\end{proof}

The argument also works in $\mathsf{ZFA}$, the set theory with atoms.
Indeed, Blass's implication and L\'evy's partition lemma hold in that
setting. In Proposition~\ref{prop:embedding}, the elements of $X$
serve only as indices for indeterminates, so the construction applies
to sets containing atoms. The subsequent arguments use finite
field extensions, sets of bases and subfields, and recursion on an
ordinal bounded by a Hartogs number, and all these constructions are
available in $\mathsf{ZFA}$.
Thus Theorem~\ref{thm:main} and
Corollary~\ref{cor:basis-principles} hold there as well.

Halbeisen~\cite[Note~94]{Halbeisen2012} asks whether, in $\mathsf{ZFA}$,
$\AC$ follows from the assertion that every vector space has a basis,
or at least from the stronger assertion that every linearly independent
subset of a vector space extends to a basis of that space. Our results
answer both questions affirmatively.

\subsection*{Acknowledgements}
This study was financed, in part, by the S\~ao Paulo Research Foundation
(FAPESP), Brasil. Process Numbers \#25/07302-0 and \#25/09425-1.

\subsection*{Declaration of AI assistance}
OpenAI's Codex, using the GPT-6 Astra model,
was used to generate novel mathematical content and assist with the writing process.
The assistance also included grammatical suggestions and bibliographical references.
All the output and suggestions were reviewed, edited and/or rewritten by the authors.